% !TeX program = xelatex
% Lesson 12: UTF-8 Chinese lecture; compile with latexmk -xelatex.
\documentclass[UTF8,11pt,a4paper]{ctexart}
\usepackage[left=21mm,right=21mm,top=20mm,bottom=22mm]{geometry}
\usepackage{amsmath,amssymb,hyperref,fvextra}
\setCJKmainfont{Noto Serif CJK SC}
\setCJKsansfont{Noto Sans CJK SC}
\setlength{\parindent}{2em}
\setlength{\parskip}{0.4em}
\hypersetup{colorlinks=true,urlcolor=blue,linkcolor=blue}
\fvset{breaklines=true,breakanywhere=true}
\begin{document}
\title{Lesson 12：GitHub 分支、PR 与自动化 CI}
\author{LLM Learn 大模型开发课程}
\date{2026 年 11 月}
\maketitle
\tableofcontents
\newpage
\section{教学目标与先修}
会使用 feature 分支、PR 和 GitHub Actions；编写仓库结构自动审计程序。\par
先修：Lesson 11。 预计 90–120 分钟。\par
\section{GitHub 协作流程}
从 main 切出 feature 分支、提交并 push，在 GitHub 打开 Pull Request；通过 Files changed 检查差异，等待 Actions 通过，再进行合并。PR 是协作审查流程而不是 Git 原生命令。

\section{CI 是什么}
CI 指提交后自动执行语法检查、单元测试、静态审查。失败时要读取作业步骤、退出码、标准输出与报错堆栈，禁止继续声称验证成功。对外部贡献的 PR 不应自动授权付费服务或密钥。

\section{仓库审计程序}
REQUIRED\_FILES 列出课程总纲、规范和工作流；find\_missing\_files 使用 Path.is\_file 返回缺失路径；main 把结构检查结果转换成进程退出码。该函数可被测试复用。

\section{练习命令}
先执行 python lessons/lesson12/code/audit\_repository.py --root .；再运行 git switch -c feature/lesson12-practice，提交文档更改并创建 PR，观察 CI 日志。

\section{核心关系或公式}
\[
\mathrm{CI}=\mathrm{syntax\ check}+\mathrm{tests}+\mathrm{artifact\ verification}
\]
\section{本课完整源码：逐行中文注释}
下列程序与本课 code 目录的实际源文件完全对应。每个模块、函数及关键可执行语句都配有中文说明。
\begin{Verbatim}[fontsize=\scriptsize,breaklines=true,breakanywhere=true]
"""Lesson 12：检查 LLM 课程仓库是否具备必要的目录和 CI 配置。"""

# 导入 argparse，允许命令行传入仓库根目录。
import argparse
# 导入 Path，进行跨平台路径检查。
from pathlib import Path

# 定义必要的仓库文件，列表本身即体现基础工程约束。
REQUIRED_FILES = (
    "README.md",
    "COURSE_PLAN.md",
    "COURSE_STANDARDS.md",
    ".github/workflows/course-code-ci.yml",
    ".github/workflows/course-latex.yml",
)


def find_missing_files(repository_root: Path) -> list[str]:
    """返回缺失文件的相对路径；没有缺失时返回空列表。"""
    # 初始化空列表，用于积累所有发现的问题。
    missing = []
    # 遍历预先定义的所有必须存在的相对路径。
    for relative_path in REQUIRED_FILES:
        # 与仓库根目录组合成完整路径。
        candidate = repository_root / relative_path
        # 若该路径不是实际存在的文件，则加入缺失清单。
        if not candidate.is_file():
            # 使用相对路径保存，避免测试日志依赖不同机器的路径。
            missing.append(relative_path)
    # 把结果返回给调用者，以便单元测试或 CLI 使用。
    return missing


def main() -> None:
    """检查指定的课程仓库，并用退出码区分合格与缺失。"""
    # 创建命令行解析器，说明输入参数的用途。
    parser = argparse.ArgumentParser(description="课程仓库基础结构审计")
    # 默认值向上寻找三层，得到本项目 Git 仓库的根目录。
    default_root = Path(__file__).resolve().parents[3]
    # 允许用户指定另外一个待检查的仓库目录。
    parser.add_argument("--root", type=Path, default=default_root)
    # 完成命令行参数解析。
    arguments = parser.parse_args()
    # 检查仓库的必要文件是否齐全。
    missing = find_missing_files(arguments.root.resolve())
    # 如果存在缺失文件，应该以失败退出码结束。
    if missing:
        # 显示每个缺失路径，便于定位问题。
        for item in missing:
            # 保持报错格式统一，供自动化日志检索。
            print("缺失:", item)
        # 退出码 1 表示审计不合格。
        parser.exit(1, "仓库结构不满足要求\n")
    # 只有全部通过时才输出成功提示。
    print("仓库核心文件检查通过。")


if __name__ == "__main__":
    # 直接执行本程序时运行主函数，导入时不会自动审计。
    main()
\end{Verbatim}
\section{练习与验收}
\begin{itemize}
\item PR 能独立创建并说明安全合并策略。
\item 能够指出 Actions 的失败步骤。
\item 审计程序对空仓库与完整仓库的测试都通过。
\end{itemize}
\section{执行记录与参考}
使用 python -m unittest discover -s tests -v 执行测试。请以本课 results/verification.md 中的真实结果和 GitHub Actions 日志为准，Mock 并非在线 API 验证。
\begin{itemize}
\item \url{https://git-scm.com/docs}
\item \url{https://docs.github.com/en/actions}
\item \url{https://developer.mozilla.org/en-US/docs/Web/HTTP}
\item \url{https://platform.openai.com/docs/api-reference/responses}
\end{itemize}
\end{document}
